% Laget automatisk av elevark/lag_elevark.py fra sim/dosert-sving.html. Endringer her blir overskrevet.
\documentclass[11pt]{article}
\usepackage{elevark}
\begin{document}
\elevarkhode{Fysikk 2 · Sirkelbevegelse}{Dosert sving}{https://fysikklab.no/sim/dosert-sving.html}{Simuleringen}

En dosert sving heller inn mot sentrum av svingen. Da peker normalkraften litt inn mot sentrum og hjelper bilen rundt. Ved én bestemt fart trengs det ingen friksjon i det hele tatt. Hva skjer hvis bilen kjører fortere eller saktere?

\section*{Oppgaver}

Start med $r = 50\,\mathrm{m}$, $θ = 15^{\circ}$, 50~km/h og tørr asfalt.

\begin{oppgaveliste}
\oppgave \fase{gjett}{Gjett}Hvilken vei peker summen av kreftene på bilen? Hvorfor er den ikke parallell med veien?
\svar
\oppgave \fase{test}{Regn}Uten friksjon: Vis at den ideelle farten er $v = \sqrt{rg \tan  θ}$. Regn den ut, og sjekk ved å velge «Uten friksjon» og finne farten der bilen holder seg på veien.
\svar
\oppgave \fase{forklar}{Forklar}Kjør i 50~km/h på tørr asfalt. Hvilken vei peker friksjonen? Hva skjer med friksjonen hvis du kjører i 30~km/h?
\svar
\oppgave \fase{test}{Undersøk}Velg is. Finn den minste og største farten bilen kan ha. Hva skjer under og over disse fartene?
\svar
\oppgave \fase{test}{Regn}Sett $θ = 0$ (flat vei) og tørr asfalt. Regn ut den største farten med $μmg = mv^{2}/r$. Hvor mye øker den største farten når veien doseres 15°?
\svar
\oppgave \fase{forklar}{Forklar}Hvorfor har velodromer og bobbaner mye brattere dosering enn vanlige veier?
\svar
\oppgave \fase{utfordring}{Utfordring}Utled uttrykket for den største farten: $v^{2} = rg (\sin  θ + μ \cos  θ) / (\cos  θ - μ \sin  θ)$. Hva skjer når $\cos  θ = μ \sin  θ$?
\svar
\end{oppgaveliste}

\end{document}
